\documentclass[12pt]{article}
\usepackage[utf8]{inputenc}
\usepackage[T1]{fontenc}
\usepackage{hyperref}
\title{Kalkoti Language Resources for K-12 Curriculum Development}
\author{Compiled by Claude Sonnet 5.5 (claude-sonnet-5-5) for Andr\'as Kornai}
\date{}
\begin{document}
\maketitle

\section{Introduction}
Kalkoti (ISO 639-3 xka; Glottocode kalk1245) is an Indo-Aryan language of Kalkot in the upper Panjkora Valley, Dir Kohistan, Khyber Pakhtunkhwa, Pakistan. Wikipedia gives about 6,000 speakers (2018). A Forum for Language Initiatives (FLI) article reports that community members claim 15,000. Henrik Liljegren describes it as at its core a Shina-type language, closely related to Palula and Sawi, with heavy convergence with a nearby Kohistani variety. Search results show a small amount of descriptive linguistics and two FLI community-development items. They show no school materials.

\section{Primary Educational Resources}
No published primers or textbooks were found. FLI reports that a workshop gave the language a written form, which is a first step towards such materials.
\begin{itemize}
\item FLI, ``Kalkoti language turns into written form'': \url{https://fli-online.org/kalkoti-language-turns-into-written-form/}
\end{itemize}

\section{Linguistic and Reference Materials}
\begin{itemize}
\item Wikipedia, Kalkoti language: \url{https://en.wikipedia.org/wiki/Kalkoti_language}
\item Numerals data sheet (Chan numerals project), Kalkoti, Dir Kohistan: \url{https://lingweb.eva.mpg.de/channumerals/Kalkoti.htm}
\item Dartmouth \emph{Linguistic Discovery} items that turned up in the search (the Liljegren work on Kalkoti is cited in the summaries, but I did not confirm these pages are that paper): \url{https://journals.dartmouth.edu/cgi-bin/WebObjects/Journals.woa/xmlpage/1/document/899} and \url{https://journals.dartmouth.edu/cgi-bin/WebObjects/Journals.woa/1/xmlpage/1/article/423}
\end{itemize}

\section{Religious and Community Resources}
\begin{itemize}
\item FLI, ``Kalkoti, a language spoken by all community children'': \url{https://fli-online.org/kalkoti-a-language-spoken-by-all-community-children/}
\end{itemize}
No scripture translations were found.

\section{Digital and Audio Resources}
None were found.

\section{Curriculum Development Context}
The language is said to be learned by almost all children in the community, so intergenerational transmission is strong, and no school use is known. FLI's written-form workshop suggests that orthography work has begun. A curriculum would need an agreed script and orthography, locally produced readers and a documented lexicon and grammar. Contact with FLI would be the natural first step, but I found no contact details.

\end{document}